% !TEX root = ../../../main.tex
\section{Transposes and linear systems}
\label{sec:linear-algebra-i:transposes-and-linear-systems}

% Source: October 8 PDF, pp. 9--12.
\subsection{Transpose}

\begin{definition}
  If \(M\) is an \(m\)-by-\(n\) matrix, its transpose \(M^T\)
  is the \(n\)-by-\(m\) matrix obtained by interchanging rows and columns,
  \[
    (M^T)^j_i=M^i_j.
  \]
\end{definition}

Transposition reverses the order of a matrix product. For
\(A\in\mathbb F^{k\times m}\) and \(B\in\mathbb F^{m\times n}\),
\[
  (AB)^T=B^TA^T\in\mathbb F^{n\times k}.
\]

\begin{theorem}[Determinant of a transpose]
  \label{thm:linear-algebra-i:determinant-transpose}
  Every square matrix satisfies \(\det(M^T)=\det M\).
\end{theorem}

\begin{proof}
  The permutation formula gives
  \begin{align*}
    \det M
      &=\sum_{\sigma\in S_n}\varepsilon(\sigma)
        M^{\sigma(1)}_1\cdots M^{\sigma(n)}_n\\
      &=\sum_{\sigma\in S_n}\varepsilon(\sigma)
        M^1_{\sigma^{-1}(1)}\cdots M^n_{\sigma^{-1}(n)}\\
      &=\sum_{\tau\in S_n}\varepsilon(\tau)
        M^1_{\tau(1)}\cdots M^n_{\tau(n)}\\
      &=\det(M^T).
  \end{align*}
  In the second line, the factors are reordered by their row indices.
  In the third, set \(\tau=\sigma^{-1}\) and use
  \(\varepsilon(\sigma^{-1})=\varepsilon(\sigma)\).
\end{proof}

% PDF p.11 supplies this factorization argument but no construction of PLU.
% Keep it conditional rather than inventing a factorization theorem or proof.
\begin{remark}
  For a factorization \(M=PLU\), where \(P\) is a permutation matrix,
  \(L\) is lower triangular, and \(U\) is upper triangular, the same
  determinant identity has the calculation
  \begin{align*}
    \det(M^T)&=\det(U^TL^TP^T)\\
      &=(\det U)(\det L)(\det P)\\
      &=\det(PLU)=\det M.
  \end{align*}
  Transposition preserves the diagonal product of a triangular matrix
  and the determinant of a permutation matrix.
\end{remark}

\subsection{Square systems of equations}

Consider \(Ax=b\) with \(A\in\R^{n\times n}\) and \(b\in\R^n\).
If \(\det A\ne0\), the matrix is invertible and the unique solution is
\[
  x=A^{-1}b.
\]
If \(\det A=0\), then
\[
  \dim\ker A>0,\qquad \dim\operatorname{im}A<n.
\]
A solution exists exactly when \(b\in\operatorname{im}A\).
If \(x_0\) is one solution, every vector \(x_0+v\), with
\(v\in\ker A\), is also a solution, since
\[
  A(x_0+v)=Ax_0+Av=b.
\]
Thus a singular square system over \(\R\) has either no solution or
infinitely many solutions. Its full solution set, when nonempty, is
\(x_0+\ker A\), as in \autoref{sec:linear-algebra-i:kernel}.
\autoref{fig:linear-algebra-i:affine-solution-fiber} illustrates this
translation when the kernel is a line in \(\R^2\).
\begin{marginfigure}
  \centering
  % A rank-one real planar schematic; the two infinite lines are parallel.
\begin{tikzpicture}[
  x=10mm,y=10mm,font=\footnotesize,line width=0.45pt,
  vector/.style={-{Latex[length=1.6mm,width=1.1mm]}},
  continued/.style={{Latex[length=1.6mm,width=1.1mm]}-{Latex[length=1.6mm,width=1.1mm]}}
]
  \coordinate (O) at (0,0);
  \coordinate (X) at (0.2,1.5);
  \coordinate (V) at (1.4,0.56);
  \coordinate (P) at ($(X)+(V)$);
  \draw[continued,black!50] (-0.7,-0.28) -- (2.8,1.12);
  \draw[continued] (-0.5,1.22) -- (3,2.62);
  \draw[vector] (O) -- (X);
  \draw[vector,dashed,black!60] (O) -- (V);
  \draw[vector] (X) -- (P);
  \fill (O) circle[radius=1.3pt];
  \fill (X) circle[radius=1.3pt];
  \fill (P) circle[radius=1.3pt];
  \node[below left=2pt] at (O) {\(0\)};
  \node[above left=3pt] at (X) {\(x_0\)};
  \node[below=4pt] at (V) {\(v\)};
  \node[above=4pt] at (P) {\(x_0+v\)};
  \node[below=4pt] at (2.4,0.96) {\(\ker A\)};
  \node at (1.2,3.3) {\(x_0+\ker A\)};
  \path (current bounding box.south) ++(0,-2mm);
\end{tikzpicture}

  \caption[A solution set as a translated kernel]{A real-plane schematic
    for a nonempty singular system. Moving from \(x_0\) by any kernel
    vector \(v\) gives another solution \(x_0+v\). The two lines are parallel.}
  \label{fig:linear-algebra-i:affine-solution-fiber}
\end{marginfigure}
